% Part: first-order-logic
% Chapter: natural-deduction
% Section: proof-theoretic-notions

\documentclass[../../../include/open-logic-section]{subfiles}

\begin{document}

\iftag{FOL}
      {\olfileid{fol}{ntd}{ptn}}
      {\olfileid{pl}{ntd}{ptn}}

\olsection{Gagasan Teoretis-Bukti}

\begin{editorial}
Bagean iki nglumpukake definisi relasi provabilitas lan konsistensi
kanggo deduksi natural.
\end{editorial}

\begin{explain}
Kaya nalika kita nemtokake sawetara gagasan semantis wigati
(validitas, akibat, lan bisa dicukupi), saiki kita nemtokake
\emph{gagasan teoretis-bukti} kang cocog. Gagasan iki ora ditemtokake
kanthi nggunakake panyukupan !!{sentence}s ing !!{structure}s, nanging
kanthi nggunakake !!{derivability} utawa !!{nonderivability}
!!{sentence}s tartamtu saka ukara-ukara liyane. Panemuan yen
gagasan-gagasan kasebut padha iku panemuan wigati. Kasunyatan manawa
gagasan kasebut padha dadi isi \emph{teorema kasahihan} lan
\emph{teorema kelengkapan}.
\end{explain}


\begin{defn}[Teorema]
!!{sentence}~$!A$ iku \emph{teorema} yen ana !!a{derivation}
saka~$!A$ ing deduksi natural kang kabeh asumsine wis
!!{discharged}. Kita nulis $\Proves !A$ yen $!A$ iku teorema, lan
$\Proves/ !A$ yen dudu teorema.
\end{defn}

\begin{defn}[!!^{derivability}]
!!^a{sentence} $!A$ \emph{bisa !!{derivable} saka} himpunan
!!{sentence}s~$\Gamma$, ditulis $\Gamma \Proves !A$, yen ana
!!{derivation} kang kesimpulane~$!A$ lan saben asumsine wis
!!{discharged} utawa kalebu ing~$\Gamma$. Yen $!A$ ora bisa
!!{derivable} saka $\Gamma$, kita nulis $\Gamma \Proves/ !A$.
\end{defn}

\begin{defn}[Konsistensi]
Himpunan !!{sentence}s~$\Gamma$ iku \emph{ora konsisten} yen lan mung
yen $\Gamma \Proves \lfalse$. Yen ora mangkono tumrap $\Gamma$, yaiku
yen $\Gamma \Proves/ \lfalse$, kita ngarani himpunan iku
\emph{konsisten}.
\end{defn}

\begin{prop}[Refleksivitas]
\ollabel{prop:reflexivity}
Yen $!A \in \Gamma$, mula $\Gamma \Proves !A$.
\end{prop}

\begin{proof}
Asumsi $!A$ dhewe iku !!a{derivation} saka~$!A$ kang saben asumsi
!!{undischarged}, yaiku $!A$, kalebu ing~$\Gamma$.
\end{proof}


\begin{prop}[Monotonisitas]
\ollabel{prop:monotonicity}
Yen $\Gamma \subseteq \Delta$ lan $\Gamma \Proves !A$, mula $\Delta
\Proves !A$.
\end{prop}

\begin{proof}
Saben !!{derivation} saka $!A$ adhedhasar $\Gamma$ uga dadi
!!a{derivation} saka $!A$ adhedhasar~$\Delta$.
\end{proof}

\begin{prop}[Transitivitas]
\ollabel{prop:transitivity}
Yen $\Gamma \Proves !A$ lan $\{!A\} \cup \Delta \Proves
!B$, mula $\Gamma \cup \Delta \Proves !B$.
\end{prop}

\begin{proof}
Yen $\Gamma \Proves !A$, ana !!a{derivation}~$\delta_0$ saka~$!A$
kanthi kabeh asumsi !!{undischarged} kalebu ing~$\Gamma$. Yen
$\{!A\} \cup \Delta \Proves !B$, mula ana
!!a{derivation}~$\delta_1$ saka~$!B$ kanthi kabeh asumsi
!!{undischarged} kalebu ing~$\{!A\} \cup \Delta$. Saiki gatekna:
\begin{prooftree}
  \AxiomC{$\Delta, \Discharge{!A}{1}$}
  \RightLabel{$\delta_1$}
  \DeduceC{$!B$}
  \DischargeRule{\Intro{\lif}}{1}
  \UnaryInfC{$!A \lif !B$}
  \AxiomC{$\Gamma$}
  \RightLabel{$\delta_0$}
  \DeduceC{$!A$}
  \RightLabel{\Elim{\lif}}
  \BinaryInfC{$!B$}
\end{prooftree}
Saiki kabeh asumsi !!{undischarged} kalebu ing $\Gamma \cup \Delta$,
mula iki nuduhake $\Gamma \cup \Delta \Proves !B$.
\end{proof}

Yen $\Gamma = \{!A_1, !A_2, \ldots, !A_k\}$ iku himpunan winates,
kita bisa nggunakake notasi ringkes $!A_1,!A_2,\ldots,!A_k \Proves
!B$ kanggo $\Gamma \Proves !B$; mligine, $!A \Proves !B$ ateges
$\{!A\} \Proves !B$.

Gatekna yen $\Gamma \Proves !A$ lan $!A \Proves !B$, mula $\Gamma
\Proves !B$. Uga ndherek yen $!A_1, \dots, !A_n \Proves !B$ lan
$\Gamma \Proves !A_i$ kanggo saben~$i$, mula $\Gamma \Proves !B$.

\begin{prop}
\ollabel{prop:incons}
Pratelan-pratelan iki ekuivalen.
    \begin{enumerate}
      \item \( \Gamma \) ora konsisten.
      \item \( \Gamma \Proves {!A} \) kanggo saben !!{sentence}~\( {!A} \).
      \item \( \Gamma \Proves {!A} \) lan \( \Gamma \Proves \lnot {!A} \) kanggo sawijining !!{sentence}~\( {!A} \).
    \end{enumerate}
\end{prop}

\begin{proof}
Gladhen.
\end{proof}

\tagprob{FOL}
\begin{prob}
Buktekna \olref[fol][ntd][ptn]{prop:incons}
\end{prob}
\tagendprob

\tagprob{notFOL}
\begin{prob}
Buktekna \olref[pl][ntd][ptn]{prop:incons}
\end{prob}
\tagendprob

\begin{prop}[Kekompakan]
  \ollabel{prop:proves-compact}
  \begin{enumerate}
  \item Yen $\Gamma \Proves !A$, mula ana subhimpunan winates
    $\Gamma_0 \subseteq \Gamma$ saengga $\Gamma_0 \Proves !A$.
  \item Yen saben subhimpunan winates saka~$\Gamma$ konsisten,
    mula $\Gamma$ konsisten.
  \end{enumerate}
\end{prop}

\begin{proof}
  \begin{enumerate}
    \item Yen $\Gamma \Proves !A$, mula ana
      !!a{derivation}~$\delta$ saka~$!A$ adhedhasar~$\Gamma$. Ayo
      $\Gamma_0$ dadi himpunan asumsi !!{undischarged} ing~$\delta$.
      Amarga saben !!{derivation} iku winates, $\Gamma_0$ mung bisa
      ngemot !!{sentence}s kanthi cacah winates. Dadi, $\delta$ iku
      !!a{derivation} saka~$!A$ adhedhasar sawijining $\Gamma_0
      \subseteq \Gamma$ kang winates.
    \item Iki kontrapositif saka (1) kanggo kasus khusus $!A \ident
      \lfalse$.
  \end{enumerate}
\end{proof}

\end{document}
